% GENERATED FILE. Edit canonical lessons, then run npm run generate:books.
% canonical-source-hash: fnv1a64:03ccdea67a18aec3
% canonical-lessons: FR-C308-naissance, FR-C308-enceinte, FR-C308-accouchement, FR-C308-allaiter, FR-C308-berceau

\chapter{Une naissance}
\label{ch:fr-une-naissance}
\hlchaptermodality{308}

\begin{chapteropening}
\textbf{By the end of this chapter:} \emph{I can talk about a birth: the birth, being pregnant, the delivery, breastfeeding and the cradle.}

Complete the last lesson of chapter 308: I can talk about a birth: the birth, being pregnant, the delivery, breastfeeding and the cradle.
\end{chapteropening}

\section[la naissance]{la naissance — the birth}

\label{lesson:FR-C308-naissance}



\begin{quote}
Before the new one: say the French for the disaster victim, then the French for the scratch.
\end{quote}

\subsection*{You'll want to know: la naissance}
\textbf{la naissance} — \textquotedblleft{}the birth\textquotedblright{}.

\begin{quote}
Nous avons fêté la naissance de notre fille.
\end{quote}

\begin{quote}
\emph{We celebrated the birth of our daughter.}
\end{quote}

\emph{La date de naissance} — date of birth; \emph{de naissance} — from birth.

\begin{cousinweb}
\emph{naissance} is built on the stem of \emph{naître} \textquotedblleft{}to be born\textquotedblright{}, which is inherited from Vulgar Latin \emph{nāscere}, from Latin \emph{nāscī} \textquotedblleft{}to be born\textquotedblright{} (Lewis \& Short).
\end{cousinweb}

\begin{grammarlens}[title={What you've built}]
One more word for telling, explaining and describing at B1.
\end{grammarlens}

\subsection*{Your turn}
\begin{itemize}
\raggedright
  \item \emph{Say it:} \emph{la naissance}
  \item \emph{Say it:} \emph{la naissance}, once more
  \item \emph{Say it:} a sentence of your own with \emph{la naissance}
  \item \emph{Recall:} say the French for the disaster victim, then the French for the scratch, then say \emph{la naissance} again
\end{itemize}

\begin{tcolorbox}[breakable,colback=teal!4,colframe=teal!35!black,title={Before you move on}]
What does la naissance mean? (\textquotedblleft{}The birth\textquotedblright{}.) Say it once more.
\end{tcolorbox}
\section[enceinte]{enceinte — pregnant}

\label{lesson:FR-C308-enceinte}



\begin{quote}
Before the new one: say the French for the scratch, then the French for the birth.
\end{quote}

\subsection*{You'll want to know: enceinte}
\textbf{enceinte} — \textquotedblleft{}pregnant\textquotedblright{}.

\begin{quote}
Ma sœur est enceinte de six mois.
\end{quote}

\begin{quote}
\emph{My sister is six months pregnant.}
\end{quote}

\emph{Enceinte de six mois} — six months pregnant. As a noun, \emph{une enceinte} is a loudspeaker or a city wall.

\begin{cousinweb}
\emph{enceinte} is inherited from Late Latin \emph{incincta} (EtymDB), usually read as \textquotedblleft{}without a belt\textquotedblright{}, from \emph{cingere} \textquotedblleft{}to gird\textquotedblright{} (Lewis \& Short).
\end{cousinweb}

\begin{grammarlens}[title={What you've built}]
One more word for telling, explaining and describing at B1.
\end{grammarlens}

\subsection*{Your turn}
\begin{itemize}
\raggedright
  \item \emph{Say it:} \emph{enceinte}
  \item \emph{Say it:} \emph{enceinte}, once more
  \item \emph{Say it:} a sentence of your own with \emph{enceinte}
  \item \emph{Recall:} say the French for the scratch, then the French for the birth, then say \emph{enceinte} again
\end{itemize}

\begin{tcolorbox}[breakable,colback=teal!4,colframe=teal!35!black,title={Before you move on}]
What does enceinte mean? (\textquotedblleft{}Pregnant\textquotedblright{}.) Say it once more.
\end{tcolorbox}
\section[l'accouchement]{l'accouchement — the delivery, giving birth}

\label{lesson:FR-C308-accouchement}



\begin{quote}
Before the new one: say the French for the birth, then the French for pregnant.
\end{quote}

\subsection*{You'll want to know: l'accouchement}
\textbf{l'accouchement} — \textquotedblleft{}the delivery, giving birth\textquotedblright{}.

\begin{quote}
L'accouchement s'est bien passé.
\end{quote}

\begin{quote}
\emph{The delivery went well.}
\end{quote}

\emph{Un accouchement à domicile} — a home birth.

\begin{cousinweb}
\emph{accouchement} is built on a verb from Old French \emph{acouchier} \textquotedblleft{}to lie down\textquotedblright{}, itself built on \emph{coucher}, which is inherited from Latin \emph{collocāre} \textquotedblleft{}to place, to lay\textquotedblright{} (Lewis \& Short).
\end{cousinweb}

\begin{grammarlens}[title={What you've built}]
One more word for telling, explaining and describing at B1.
\end{grammarlens}

\subsection*{Your turn}
\begin{itemize}
\raggedright
  \item \emph{Say it:} \emph{l'accouchement}
  \item \emph{Say it:} \emph{l'accouchement}, once more
  \item \emph{Say it:} a sentence of your own with \emph{l'accouchement}
  \item \emph{Recall:} say the French for the birth, then the French for pregnant, then say \emph{l'accouchement} again
\end{itemize}

\begin{tcolorbox}[breakable,colback=teal!4,colframe=teal!35!black,title={Before you move on}]
What does l'accouchement mean? (\textquotedblleft{}The delivery, giving birth\textquotedblright{}.) Say it once more.
\end{tcolorbox}
\section[allaiter]{allaiter — to breastfeed}

\label{lesson:FR-C308-allaiter}



\begin{quote}
Before the new one: say the French for pregnant, then the French for the delivery.
\end{quote}

\subsection*{You'll want to know: allaiter}
\textbf{allaiter} — \textquotedblleft{}to breastfeed\textquotedblright{}.

\begin{quote}
Elle a allaité son bébé pendant six mois.
\end{quote}

\begin{quote}
\emph{She breastfed her baby for six months.}
\end{quote}

Built on \emph{lait}, milk.

\begin{cousinweb}
\emph{allaiter} is usually traced to Late Latin \emph{allactāre} \textquotedblleft{}to suckle\textquotedblright{}, built on \emph{lac}, \emph{lactis} \textquotedblleft{}milk\textquotedblright{} (Lewis \& Short); EtymDB lists its Catalan and Italian cousins.
\end{cousinweb}

\begin{grammarlens}[title={What you've built}]
One more word for telling, explaining and describing at B1.
\end{grammarlens}

\subsection*{Your turn}
\begin{itemize}
\raggedright
  \item \emph{Say it:} \emph{allaiter}
  \item \emph{Say it:} \emph{allaiter}, once more
  \item \emph{Say it:} a sentence of your own with \emph{allaiter}
  \item \emph{Recall:} say the French for pregnant, then the French for the delivery, then say \emph{allaiter} again
\end{itemize}

\begin{tcolorbox}[breakable,colback=teal!4,colframe=teal!35!black,title={Before you move on}]
What does allaiter mean? (\textquotedblleft{}To breastfeed\textquotedblright{}.) Say it once more.
\end{tcolorbox}
\section[le berceau]{le berceau — the cradle}

\label{lesson:FR-C308-berceau}



\begin{quote}
Before the new one: say the French for the delivery, then the French for to breastfeed.
\end{quote}

\subsection*{You'll want to know: le berceau}
\textbf{le berceau} — \textquotedblleft{}the cradle\textquotedblright{}.

\begin{quote}
Le bébé dort dans son berceau.
\end{quote}

\begin{quote}
\emph{The baby is sleeping in its cradle.}
\end{quote}

Figurative: \emph{le berceau d'une civilisation}.

\begin{cousinweb}
\emph{berceau} comes from Old French \emph{berçuel}, inherited from Vulgar Latin \emph{berciolum} \textquotedblleft{}cradle\textquotedblright{} (EtymDB), which is usually traced to Gaulish.
\end{cousinweb}

\begin{grammarlens}[title={What you've built}]
One more word for telling, explaining and describing at B1.
\end{grammarlens}

\subsection*{Your turn}
\begin{itemize}
\raggedright
  \item \emph{Say it:} \emph{le berceau}
  \item \emph{Say it:} \emph{le berceau}, once more
  \item \emph{Say it:} a sentence of your own with \emph{le berceau}
  \item \emph{Recall:} say the French for the delivery, then the French for to breastfeed, then say \emph{le berceau} again
\end{itemize}

\begin{tcolorbox}[breakable,colback=teal!4,colframe=teal!35!black,title={Before you move on}]
What does le berceau mean? (\textquotedblleft{}The cradle\textquotedblright{}.) Say it once more.
\end{tcolorbox}
